% PlanetMath proposal draft for arXiv:2602.05990
% Source corpora: arXiv math.CT eprints, storage/mo-processed-gpu, storage/math-processed-gpu
\section*{Categories graded by group homomorphisms}
\subsection*{Source}
\begin{itemize}
  \item arXiv:2602.05990, Jonathan Davies, ``Categories graded by group homomorphisms'' (2026-02-11).
  \item MathOverflow 401964 on graded categories and non-abelian cohomology.
  \item Math.SE 198017 on graded versions of Baer's criterion for injectivity.
\end{itemize}
\subsection*{Synopsis}
Davies extends the Sözer--Virelizier framework of $\chi$-graded bicategories to general homomorphisms $\tau\colon H\to G$. A $\tau$-graded category carries degrees on both objects ($G$) and morphisms ($H$), interpolation between enriched and graded settings. Main points:
\begin{enumerate}
  \item Definition of the $2$-category $\Cat_{\tau}$ of $\tau$-graded categories, functors, and transformations, together with ``half-enriched'' Yoneda embeddings.
  \item Proof of a Yoneda lemma in this setting: representables factor through $\tau$-graded functors with prescribed shifts. This addresses MO concerns about outer actions controlling graded equivalences.
  \item Construction of shift objects and the resulting description of semisimple $\tau$-graded categories as extensions of indecomposables indexed by subgroup data in $\ker(\tau)$. The semisimple classification resembles graded Morita theory on Math.SE 198017.
  \item Comparison with $\tau$-module categories, module categories over groupoid algebras induced by $H \to G$, and the resulting $2$-equivalence $\Cat_{\tau}^{\shifts} \simeq \ModCat[H]_{\tau}$.
\end{enumerate}
\subsection*{MathOverflow cues}
\begin{description}
  \item[MO 401964] asks how cohomological obstructions classify $G$-graded categories with outer actions; Davies provides explicit data via $\tau$-graded shift objects and obstruction classes in $H^2(H,\Hom(G,\Bbbk^{\times}))$.
\end{description}
\subsection*{Math.SE anchors}
\begin{description}
  \item[Math.SE 198017] discusses graded analogues of Baer's criterion. The paper’s half-enriched Yoneda lemma supplies the needed homological control on $\tau$-graded injectives, linking local grading data with finitely presented morphisms.
\end{description}
\subsection*{Encyclopedia outline}
\begin{enumerate}
  \item \textbf{Definition.} Give the precise data of a $\tau$-graded category, functor, and natural transformation; illustrate with $\tau=\operatorname{id}_G$ and $\tau$ injective/surjective extremes.
  \item \textbf{Half-enriched Yoneda lemma.} Present the statement, explain how representables carry degree shifts indexed by $H$, and sketch proof strategies (use of $\tau$-decorated oplax limits).
  \item \textbf{Shift objects and semisimplicity.} Summarize the $\tau$-shift construction and the classification of semisimple $\tau$-graded categories via subgroups of $\ker(\tau)$ and $2$-cocycles.
  \item \textbf{$\tau$-module categories.} Describe $\ModCat[H]_{\tau}$ and the $2$-equivalence to $\Cat_{\tau}^{\shifts}$; relate to module categories over skew-group algebras.
  \item \textbf{Examples.} Work out examples for crossed $\chi$-graded categories, finite $G$-sets with $H$-labelled morphisms, and $\tau$ induced by group extensions.
\end{enumerate}
\subsection*{Action items}
\begin{itemize}
  \item Cross-check existing PlanetMath entries on graded categories, crossed modules, and enriched categories to avoid duplication.
  \item Prepare illustrative diagrams showing object/morphism degrees and the effect of $\tau$.
  \item Summarize the Yoneda proof and highlight differences with enriched Yoneda lemmas.
  \item Draft example computations for $\tau$ arising from $1 \to H \to \widetilde{G} \to G \to 1$ extensions.
\end{itemize}
